\KoreanAbstractHeader

휴대용 기기에서 대형 언어 모델을 실행할 때에는 DRAM에 올리지 못한 가중치를 플래시에 두고, 입력마다 중요한 행만 선택하여 읽는 방법이 쓰인다. SSD의 처리량은 요청 크기가 커질수록 증가하다가 포화되므로, 같은 수의 행을 읽더라도 선택한 행이 얼마나 연속되어 있는지에 따라 읽기 시간이 달라진다. Yang 등의 Neuron Chunking은 이 점을 반영하여, 추정 읽기 시간당 보존 중요도를 최대화하는 행 선택 문제를 제안하였다.

본 논문은 이 문제가 행 수의 상한만 제약하기 때문에 단일 연속 구간이 항상 최적해가 됨을 보이고, 정확히 $R$개의 행을 선택하는 문제로 바꾼다. 이어서 처리량 포화를 두 직선으로 근사한 two-line 지연 모델을 도입한다. 이 모델에서 읽기 비용은 행 수에 비례하는 비용과 짧은 구간의 추가 비용으로 분해된다. 이를 이용하여 동적계획법과 Dinkelbach 방법으로 최적 mask를 구하는 Global DP와, 일정 길이의 타일을 단위로 행을 선택하는 Tiles를 제시한다. 또한 모든 구간의 길이가 전환점 이상이면, Top-$R$ 대비 중요도 보존율이 최적 비율 대비 달성률의 하한이 됨을 보인다.

노트북의 NVMe SSD에서 측정한 읽기 비용에 대하여 two-line 모델의 test MAPE는 3.4\%였다. 언어 모델에서 수집한 중요도로 비교한 결과, $896\times896$ 가중치 행렬의 50\% 예산에서 Tiles는 최적 비율의 95.9\%에, Neuron Chunking의 greedy 방법은 96.8\%에 도달하였다. 실제 측정에서 두 방법의 SSD 읽기 시간은 비슷하였으나, 행을 선택하는 시간은 Tiles가 약 100분의 1이었다. 그 결과 선택과 읽기를 합한 전체 시간은 Tiles가 약 3분의 1이었고, 보존 중요도는 3.6\%p 낮았다.

\KoreanKeywords{선택적 가중치 로딩, SSD 처리량 포화, activation sparsification, 동적계획법, 분수 계획법, on-device 추론}
